\documentstyle[% 


righttag% 


,11pt% 


,amssymb% 


,russian%               remove in Eng. 


% ,izvru%                 Set izven% in Eng. 


,izvru-ks%             for brief communications 


]{amsart} 


 


% 


 


\newcommand{\eqdef}{\overset{\mathrm{def}}{=}} 


\newcommand{\up}[2]{\overset{#1}{#2}{}} 


\renewcommand{\baselinestretch}{0.96}


 


%\hoffset=-15mm%                  For Eng. 


\setcounter{page}{74} 


\god{2005} 


\nomer{11~(522)} %                        Remove in Eng. 


%\nomer{Vol.\,49, No.\,11}%               Set in Eng. 


%\str{\pageref{begin}--\pageref{end}}%  Set in Eng. 


\udk{514.756} 


 


\author{\it Т.Н.\,Андреева} 


% NB:   ^^ remove in Eng. 


\title{Линейные связности, 


индуцируемые полным оснащением гиперповерхности 


конформного пространства} 


 


\date{} 


%\pagestyle{myheadings}%         Set in Eng. 


\pagestyle{plain} %              Remove in Eng. 


\begin{document} 


%\thispagestyle{plain}%          Set in Eng. 


\maketitle 


\label{begin} 


 


 


{\bf 1.} Рассмотрим гиперповерхность [1] $V_{n-1}\subset C_n$ ($n\ge3$); 


отнесем ее к 


полуизотропному [2] полуортогональному ($g_{in}=(A_i A_n)=0$) 


($i,j,k,l,s,t=\ov{1,n-1}$) 


реперу $R=\{A_\lambda\}$ ($\lambda,\mu,\rho=\ov{0,n+1}$) первого порядка. 


Уравнения инфинитезимального  


перемещения репера $R$ имеют вид $dA_\lambda=\omega^\mu_\lambda A_\mu$, 


где дифференциальные  


формы Пфаффа $\omega^\mu_\lambda$ удовлетворяют уравнениям структуры 


конформного пространства $C_n$ [3], [4]: 


$D\omega^\mu_\lambda=\omega^\rho_\lambda\wedge\omega^\mu_\rho$, а также 


линейным 


зависимостям 


\begin{gather*} 


\omega^{n+1}_0=\omega^0_{n+1}=0,\ \ \omega^0_0+\omega^{n+1}_{n+1} 


=0,\ \ \omega^0_I+g_{IK}\omega^K_{n+1}=0,\\ 


\omega^{n+1}_I+g_{IK}\omega^K_0=0,\ \  


dg_{IL}-g_{IK}\omega^K_L-g_{KL}\omega^K_I=0 


\quad (I,J,K,L=\ov{1,n}). 


\end{gather*} 


 


Если скалярные произведения $(A_\lambda A_\mu)$ элементов выбранного 


репера 


обозначить через $g_{\lambda\mu}$, то [3], [4] 


$$ 


\|g_{\lambda\mu}\|= 


\begin{Vmatrix} 


0 & 0 & 0 & 1 \\ 0 & g_{ij} & 0 & 0 \\ 0 & 0 & g_{nn} & 0 \\ 


1 & 0 & 0 & 0 


\end{Vmatrix} 


,\quad g_{\lambda\mu}=g_{\mu\lambda}, 


$$ 


причем метрический тензор $g_{ij}$ и относительный инвариант $g_{nn}$ 


являются  


невырожденными: 


\begin{gather*} 


g_{il}g^{lj}=\delta^j_i,\ \ g_{nn}g^{nn}=1,\ \  


dg_{ij}-g_{ik}\omega^k_j-g_{kj}\omega^k_i=0,\\ 


dg^{ij}+g^{ik}\omega^j_k+g^{kj}\omega^i_k=0,\ \  


d\ln g_{nn}-2\omega^n_n=0,\ \ d\ln g^{nn} 


+2\omega^n_n=0. 


\end{gather*} 


 


В этом репере справедливо 


$$ 


\omega^n_0=0,\ \ \omega^n_i=\Lambda^n_{ij}\omega^j_0,\ \  


\omega^i_n=\Lambda^i_{nt}\omega^t_0,\ \  


\Lambda^n_{[ij]}=0,\ \ g_{ij}\Lambda^j_{nk} 


+g_{nn}\Lambda^n_{ik}=0. 


$$ 


 


{\bf 2.} Предположим, что задано касательное оснащение [1], [5]  


гиперповерхности $V_{n-1}\subset C_n$ полем гиперсфер $P_n=A_n+x^0_n 


A_0$, определяемым полем квазитензора $x^0_n$ [1]:  


\begin{equation} 


dx^0_n+x^0_n(\omega^0_0-\omega^n_n)+\omega^0_n=x^0_{nk}\omega^k_0. 


\end{equation} 


 


Охваты $\Lambda_n\eqdef-\frac{1}{n-1}\Lambda^j_{nj}$, 


$\wt\Lambda_n\eqdef\frac{1}{n-1}g^{jk}g_{nn}\Lambda^n_{kj}$ удовлетворяют 


уравнению (1); 


кроме того, данные охваты совпадают 


$\Lambda_n=\wt\Lambda_n=-\frac{1}{n-1}\Lambda^j_{nj}$. Таким образом, во 


второй дифференциальной окрестности охват $\Lambda_n$ внутренним образом  


определяет касательное оснащение гиперповерхности $V_{n-1}\subset C_n$. 


 


Возьмем систему из $(n+1)^2$ форм Пфаффа $\{\up{0}{\Omega}^a_b\}$ 


($a,b,c=\ov{0,n-1};n+1$): 


\begin{equation} 


\begin{array}{c} 


\up{0}{\Omega}^j_0=\omega^j_0,\ \ \up{0}{\Omega}^0_0= 


\omega^0_0,\ \ \up{0}{\Omega}^0_0+\up{0}{\Omega}^{n+1}_{n+1}=0, 


\ \ \up{0}{\Omega}^0_i=\omega^0_i-x^0_n\omega^n_i- 


\dfrac{1}{2}g^{nn}(x^0_n)^2\omega^{n+1}_i,\\[3mm] 


\up{0}{\Omega}^j_i=\omega^j_i,\ \  


\up{0}{\Omega}^j_{n+1}=-g^{jk}\up{0}{\Omega}^0_k,\ \  


\up{0}{\Omega}^{n+1}_i=\omega^{n+1}_i,\ \  


\up{0}{\Omega}^{n+1}_0=\omega^{n+1}_0=0,\ \  


\up{0}{\Omega}^0_{n+1}=\omega^0_{n+1}=0. 


\end{array} 


\end{equation} 


Согласно работе [5] система форм (2) удовлетворяет структурным 


уравнениям пространства конформной связности $\up{0}{C}_{n-1,n-1}$ с 


$(n-1)$-мерной базой $V_{n-1}$ и $(n-1)$-мерными слоями, являющимися 


конформными пространствами 


$C_{n-1}(u^i)\equiv C_{n-1}(u)$: 


\begin{equation} 


D\omega^i_0=\omega^k_0\wedge(\omega^i_k- 


\delta^i_k\omega^0_0),\quad D\up{0}{\Omega}^a_b= 


\up{0}{\Omega}^c_b\wedge\up{0}{\Omega}^a_c+ 


\frac{1}{2}\up{0}{R}^a_{bkl}\omega^k_0\wedge\omega^l_0, 


\end{equation} 


где $d\up{0}{R}^a_{bks}+2\up{0}{R}^a_{bks} \up{0}{\Omega}^0_0- 


\up{0}{R}^a_{bkl}\up{0}{\Omega}^l_s- \up{0}{R}^a_{bls}\up{0}{\Omega}^l_k- 


\up{0}{R}^a_{cks} \up{0}{\Omega}^c_b+ \up{0}{R}^c_{bks} 


\up{0}{\Omega}^a_c= \up{0}{R}^a_{bksl}\omega^l_0$. 


В структурных уравнениях (3) компоненты тензора кривизны-кручения 


$\up{0}{R}^a_{bkl}$ пространства $\up{0}{C}_{n-1,n-1}$ имеют следующие 


строения: 


\begin{equation} 


\begin{array}{c} 


\up{0}{R}^0_{0kl}=\up{0}{R}^j_{0kl}= 


\up{0}{R}^{n+1}_{0kl}=\up{0}{R}^0_{n+1\,kl}= 


\up{0}{R}^{n+1}_{ikl}= 


\up{0}{R}^{n+1}_{n+1\,kl}=0,\quad 


\up{0}{R}^j_{ikl}=2(\Lambda{}^n_{i[k}-g^{nn}x^0_n 


g{}_{i[k})(\Lambda{}^j_{|n|l]}+\delta{}^j_{l]}x^0_n),\\[2mm] 


\up{0}{R}^j_{n+1\,kl}=-g^{js}\up{0}{R}^0_{skl},\quad 


\up{0}{R}^0_{ikl}=2[g^{nn}x^0_n x{}^0_{n[k}g{}_{l]i}- 


x{}^0_{n[k}\Lambda{}^n_{l]i}]. 


\end{array} 


\end{equation} 


 


Каждая из систем функций $a^n_{ik}(x)\eqdef\Lambda^n_{ik}- g^{nn}x^0_n 


g_{ik}$, $\wt\Lambda{}^j_{nl}(x)\eqdef \Lambda^j_{nl}+\delta^j_l x^0_n$ 


есть тензор порядка не ниже второго. Тогда компоненты тензора 


$\up{0}{R}^j_{ikl}$ в формулах (4) 


можно записать в виде $\up{0}{R}^j_{ikl}=2a{}^n_{i[k}(x) 


\wt\Lambda{}^j_{|n|l]}(x)$. 


 


Для пространства $\up{0}{C}_{n-1,n-1}$ без кручения 


($\up{0}{R}^j_{0kl}=0$) выполняются аналоги 


тождеств Риччи 


\begin{equation} 


\up{0}{R}^i_{(kst)}=\up{0}{R}^0_{(kst)}= 


\up{0}{R}{}^l_{i(ks} g{}^{\phantom{l}}_{t)l} =0; 


\end{equation} 


кроме того, каждая из систем функций $\{\up{0}{R}^i_{kst}\}$, 


$\{\up{0}{R}^i_{kst},\up{0}{R}^0_{kst}\}$, 


$\{\up{0}{R}^i_{kst},\up{0}{R}^i_{n+1\,st}\}$ образует 


тензор. 


 


 


\begin{thm} 


Инвариантное касательное оснащение гиперповерхности 


$V_{n-1}\subset C_n$ полем гиперсфер $P_n$ индуцирует пространство 


конформной  


связности $\up{0}{C}_{n-1,n-1}$ с полем метрического тензора $g_{ij}$, 


определяемое системой 


$(n+1)^2$ форм Пфаффа $(2)$, причем это пространство является 


пространством 


без кручения $(\up{0}{R}^j_{0kl}=0);$ компоненты тензора 


кривизны-кручения пространства 


$\up{0}{C}_{n-1,n-1}$ имеют строения $(4)$, причем выполняются аналоги 


тождеств Риччи $(5)$. 


\end{thm} 


 


{\bf 3.} Для того чтобы касательная гиперсфера $P_n=A_n+x^0_nA_0$ была 


неподвижна, необходимо и достаточно выполнение одной из двух групп 


соотношений 


\begin{equation} 


\wt \Lambda{}^l_{nk}(x)\eqdef \Lambda^l_{nk}+ 


\delta^l_k x^0_n=0,\quad x^0_{nk}=0. 


\end{equation} 


В случае неподвижности касательной гиперсферы $P_n$ функция $x^0_n$ в 


силу (6) совпадает с\linebreak $\Lambda_n$: $x^0_n\equiv\Lambda_n$, а исходная 


гиперповерхность $V_{n-1}\subset C_n$ совпадает с 


неподвижной гиперсферой $P_n=A_n+\Lambda_nA_0$. 


Кроме того, неподвижность касательной гиперсферы $P_n$ равносильна 


тому, что пространство конформной связности $\up{0}{C}_{n-1,n-1}$ 


является плоским. 


 


 


\begin{thm} 


Поле касательных гиперсфер $P_n$ гиперповерхности $V_{n-1}\subset C_n$ 


вырождается в одну гиперсферу тогда и только тогда, когда справедлива 


одна из двух эквивалентных групп соотношений $(6)$, или, что то же самое, 


когда индуцируемое пространство конформной связности 


$\up{0}{C}_{n-1,n-1}$ является  


плоским\/$;$ при этом исходная гиперповерхность $V_{n-1}$ совпадает с 


неподвижной 


гиперсферой $P_n=A_n+\Lambda_nA_0$. 


\end{thm} 


 


Гиперповерхность $V_{n-1}\subset C_n$, касательно оснащенная неподвижной 


гиперсферой $P_n$, существует с произволом $n+1$ произвольных постоянных, 


которые определяют координаты центра неподвижной гиперсферы 


$V_{n-1}\subset C_n$ и 


ее радиус. 


\medskip 


 


{\bf 4.} Предположим, что задано полное оснащение [1], [5] 


гиперповерхности 


$V_{n-1}\subset C_n$, т.\,е. кроме ее касательного оснащения полем 


гиперсфер $P_n$, 


определяемым полем квазитензора $x^0_n$, задано нормальное оснащение 


подмногообразия $V_{n-1}$ полем окружностей $[P_i]$, $P_i=A_i+x^0_iA_0$, 


определяемым 


полем квазитензора $x^0_i$ [1]: $dx^0_i+x^0_i\omega^0_0- x^0_j\omega^j_i+ 


\omega^0_i=x^0_{ij}\omega^j_0$. 


 


Задание поля квазитензора $x^0_i$ определяет нормализацию [6] 


пространства 


конформной связности $\up{0}{C}_{n-1,n-1}$. Справедлива 


 


 


\begin{thm} 


Инвариантное полное оснащение гиперповерхности $V_{n-1}\subset C_n$ 


полем квазитензора $x^0_l$ индуцирует нормализованное пространство 


конформной связности $\up{0}{C}_{n-1,n-1}$, определяемое полем 


окружностей $[P_i]$, 


задаваемых полем квазитензора $x^0_i$. 


\end{thm} 


 


Система функций $a^0_{ik}\eqdef\wt x{}^0_{ik}- x^0_i x^0_k+\frac{1}{2} 


g_{ik}g^{jl}x^0_j x^0_l$ образует тензор, вообще 


говоря, несимметричный $da^0_{ik}+2a^0_{ik} \up{0}{\Omega}^0_0 


-a^0_{lk}\up{0}{\Omega}^l_i- a^0_{il}\up{0}{\Omega}^l_k= 


a^0_{ikl}\up{0}{\Omega}^l_0$. Следуя работе [7], 


тензор $a^0_{ik}$ назовем основным тензором нормализации пространства 


$\up{0}{C}_{n-1,n-1}$, задаваемой


полем квазитензора $x^0_i$. В случае симметрии основного тензора 


$a^0_{ik}$ 


нормализацию пространства конформной связности $\up{0}{C}_{n-1,n-1}$ по 


аналогии с 


нормализацией проективного пространства [6] назовем 


гармонической. 


\medskip 


 


{\bf 5.} Охват $A^0_{ijk}\eqdef a^0_{ijk}- a^0_{ik}x^0_j- a^0_{kj}x^0_i+ 


(g_{ik}a^0_{lj}+g_{jk}a^0_{il})g^{lt}x^0_t-2a^0_{ij}x^0_k$ 


является тензором. 


 


Предположим, что основной тензор нормализации $a^0_{ik}$ невырожден: 


$a\eqdef|a^0_{ik}|\ne0$, $a^0_{ik}a^{kj}_0=a^0_{ki}a^{jk}_0= \delta^j_i$. 


В этом случае будем говорить, что нормализация 


пространства $\up{0}{C}{}_{n-1,n-1}$ и полное оснащение гиперповерхности 


$V_{n-1}\subset C_n$ 


являются невырожденными. 


 


Возьмем новую систему из $(n+1)^2$ форм Пфаффа 


\begin{equation} 


\ov\Omega{}^a_b=\up{0}{\Omega}^a_b+\Pi^a_{bk} 


\up{0}{\Omega}^k_0; 


\end{equation} 


потребуем, чтобы система форм $\{\ov\Omega{}^a_b\}$ удовлетворяла 


структурным уравнениям 


пространства конформной связности $\ov C_{n-1,n-1}$: 


$D\ov\Omega{}^a_b=\ov\Omega{}^c_b\wedge\ov\Omega{}^a_c+ \frac{1}{2}\ov 


R^a_{blt}\up{0}{\Omega}^l_0\wedge\up{0}{\Omega}^t_0$. 


Это требование равносильно тому, что система функций $\Pi^a_{bk}$ есть 


тензор. 


 


На формы системы $\{\ov\Omega{}^a_b\}$ 


наложим дополнительные условия, а именно, 


чтобы при преобразованиях (7) 


\begin{itemize} 


\item[1)] формы $\ov\Omega{}^k_0$, $\ov\Omega{}^0_0$, 


$\ov\Omega^{n+1}_{n+1}$ оставались без изменения, т.\,е. 


$$ 


\ov\Omega{}^k_0=\up{0}{\Omega}^k_0,\quad 


\ov\Omega{}^0_0=\up{0}{\Omega}^0_0,\quad 


\ov\Omega{}^{n+1}_{n+1}=\up{0}{\Omega}^{n+1}_{n+1}; 


$$ 


\item[2)] для пространства $\ov C_{n-1,n-1}$ метрическим тензором был 


метрический 


тензор $g_{ij}$ пространства $\up{0}{C}_{n-1,n-1}$; 


\item[3)] формы $\ov\Omega{}^a_b$ удовлетворяли соотношениям вида (2). 


\end{itemize} 


 


С учетом требований 1)--3) в качестве тензора $\Pi^j_{ik}$ (в 


предположении, что 


задана невырожденная нормализация пространства $\up{0}{C}_{n-1,n-1}$) 


можно взять охват 


$A^j_{ik}\eqdef a^{jt}_0A^0_{tik}-g^{jt}g_{si}a^{sl}_0A^0_{ltk}$: 


$\Pi^j_{ik}=A^j_{ik}$; при этом соответствующие формы $\ov\Omega{}^a_b$ 


связности обозначим через $\up{1}{\Omega}^a_b$, а само пространство 


конформной связности 


--- через $\up{1}{C}_{n-1,n-1}$. 


 


 


\begin{thm} 


Если полное оснащение гиперповерхности $V_{n-1}\subset C_n$ является 


невырожденным, то индуцируется второе пространство конформной 


связности $\up{1}{C}_{n-1,n-1}$ $($с кривизной и кручением\/$)$, 


метрический тензор которого 


совпадает с метрическим тензором $g_{ij}$ пространства 


$\up{0}{C}_{n-1,n-1}$, формами 


связности его являются формы $\up{1}{\Omega}{}^a_b$. 


\end{thm} 


 


{\bf 6.} Поле квазитензора $x^0_i$, определяющее нормализацию пространства  


конформной связности $\up{0}{C}_{n-1,n-1}$, задает нормализацию и 


индуцированного  


пространства конформной связности $\up{1}{C}_{n-1,n-1}$; при этом поля 


основных тензоров 


$a^0_{ik}$ и $\up{1}{a}^0_{ik}$ нормализованных пространств 


$\up{0}{C}_{n-1,n-1}$ и $\up{1}{C}_{n-1,n-1}$ совпадают: $a^0_{ik}\equiv 


\up{1}{a}^0_{ik}$. 


 


Каждая из систем форм $\{\up{0}{\theta}^j_0= 


\up{0}{\Omega}^j_0$, $\up{0}{\theta}^j_i=\up{0}{\Omega}^j_i- 


\delta^j_i(\up{0}{\Omega}^0_0-x^0_k\up{0}{\Omega}^k_0)+ 


g^{jl}x^0_l\up{0}{\Omega}^{n+1}_i+x^0_i\up{0}{\Omega}^j_0\}$,  


$\{\up{1}{\theta}^j_0=\up{1}{\Omega}^j_0$, 


$\up{1}{\theta}^j_i=\up{1}{\Omega}^j_i- 


\delta^j_i(\up{1}{\Omega}^0_0-x^0_k\up{1}{\Omega}^k_0)+ 


g^{jl}x^0_l\up{1}{\Omega}^{n+1}_i+x^0_i\up{1}{\Omega}^j_0\}$ удовлетворяет 


структурным 


уравнениям Картана--Лаптева [8], [9], следовательно,  


определяет пространства аффинной связности соответственно 


$\up{0}{A}_{n-1,n-1}$ и $\up{1}{A}_{n-1,n-1}$. 


 


 


\begin{thm} 


Невырожденное полное оснащение гиперповерхности $V_{n-1}\subset C_n$ 


определяет одновременную нормализацию пространств конформной 


связности $\up{0}{C}_{n-1,n-1}$ и $\up{1}{C}_{n-1,n-1}$ полем 


квазитензора $x^0_i$, при этом поля основных 


тензоров $a^0_{ik}$ и $\up{1}{a}^0_{ik}$ этих пространств совпадают\/$;$ 


аффинные связности $\up{0}{\nabla}$ и 


$\up{1}{\nabla}$, индуцируемые при этом оснащении и определяемые 


системами форм 


соответственно $\{\up{0}{\theta}^j_0, \up{0}{\theta}^j_i\}$ и 


$\{\up{1}{\theta}^j_0, \up{1}{\theta}^j_i\}$, являются вейлевыми с полем 


метрического 


тензора $g_{ik}$, причем связность $\up{0}{\nabla}$ без кручения, 


тензор кручения 


связности $\up{1}{\nabla}$ совпадает с тензором кручения пространства 


$\up{1}{C}_{n-1,n-1}$ и, вообще говоря, 


является ненулевым. Связность $\up{0}{\nabla}$ риманова тогда и только 


тогда, когда 


нормализация пространства $\up{0}{C}_{n-1,n-1}$ является гармонической. 


\end{thm} 


 


 


{\bf Замечание.} Невырожденное полное оснащение гиперповерхности 


$V_{n-1}\subset C_n$ полями функций $x^0_n$ и квазитензора $x^0_i$ в 


общем случае индуцирует 


бесчисленное множество пространств конформной связности 


$\up{p}{C}_{n-1,n-1}$ и 


аффинной связности $\up{p}{A}_{n-1,n-1}$ ($p=0,1,2,\dotsc$), 


удовлетворяющих теоремам 


соответственно 4 и 5. 


 


 


\begin{thebibliography}{9} 


 


\bibitem{1} 


Акивис М.А. {\em Инвариантное построение геометрии гиперповерхности 


конформного пространства} // Матем. сб. -- 1952. -- Т.\,31. 


-- \No\,1. -- С.\,43--75. 


\bibitem{2} 


Бушманова Г.В., Норден А.П. {\em Элементы конформной геометрии}. -- 


Казань: Изд-во Казанск. ун-та, 1972. -- 178~с. 


\bibitem{3} 


Акивис М.А. {\em К конформно-дифференциальной геометрии многомерных 


поверхностей} // Матем. сб. -- 1961. -- Т.\,53. -- \No\,1. -- С.\,53--72. 


\bibitem{4} 


Akivis M.A., Goldberg V.V. {\em Conformal differential geometry and its 


generalizations}. -- USA, 1996. -- 384~p. 


\bibitem{5} 


Столяров А.В. {\em Линейные связности на распределениях конформного 


пространства} // Изв. вузов. Математика. -- 2001. -- \No\,3. -- 


С.\,60--72. 


\bibitem{6} 


Норден А.П. {\em Пространства аффинной связности}. -- М.: Наука, 


1976. -- 432~с. 


\bibitem{7} 


Столяров А.В. {\em Внутренняя геометрия нормализованного конформного 


пространства} // Изв. вузов. Математика. -- 2002. -- \No\,11. 


-- С.\,61--70. 


\bibitem{8} 


Лаптев Г.Ф. {\em Дифференциальная геометрия погруженных многообразий. 


Теоретико-групповой метод дифференциально-геометрических исследований}  


// Тр. Моск. матем. о-ва. -- 1953. -- \No\,2. -- С.\,275--382. 


\bibitem{9} 


Евтушик Л.Е., Лумисте Ю.Г., Остиану Н.М., Широков А.П. 


{\em Дифференциально-геометрические структуры на многообразиях} // 


Итоги науки и техн. Проблемы геометрии. -- М.: ВИНИТИ, 


1979. -- Т.\,9. -- 246~с. 


 


\end{thebibliography} 


 


\vskip 0.5cm 


 


\postup{\it Чувашский государственный}{\qquad \it Поступили} 


\postup{\it педагогический университет}{\it полный текст $27.01.2005$} 


\postup{}{\it краткое сообщение $23.05.2005$} 


 


\label{end} 


\end{document} 


